\documentclass[10pt,a4paper]{amsart}
\usepackage[margin=19mm]{geometry}
\usepackage{amsmath,amssymb,mathtools}
\usepackage[scr=rsfs,cal=euler]{mathalfa}
\usepackage{xfrac}
\usepackage[unicode,hidelinks,bookmarks=false]{hyperref}
\newcommand{\ie}{i.e.\ }
\newcommand{\cf}{cf.\ }
\newcommand{\resp}{respectively\ }
\newcommand{\Subsection}{\subsection}
\providecommand{\nobreakdash}{\nobreak\hbox{-}}
\setlength{\emergencystretch}{2em}
\multlinegap0pt
\usepackage{tikz}
\usepackage{xcolor}
\usepackage{amssymb}
\newtheorem{corollary}{Corollary}
\title{Diameter-free reverse inequalities and superorthogonality}
\author{Adam Cushman, Ciprian Demeter, Shukun Wu}
\date{Source: arXiv:2609.13105v1 (CC BY 4.0)}
\begin{document}
\maketitle

\noindent\emph{Source and licence.} Diameter-free reverse inequalities and superorthogonality. Version arXiv:2609.13105v1 (CC BY 4.0), distributed by arXiv under the Creative Commons Attribution 4.0 International licence, http://creativecommons.org/licenses/by/4.0/, as recorded in the arXiv licence metadata for this exact version and shown on the abstract page https://arxiv.org/abs/2609.13105v1. Full licence notice: \texttt{licenses/arxiv-2609.13105-CC-BY-4.0.txt}.\par\smallskip

\noindent\textit{Editorial scope.} Counted unit: the article's corollary cor:exact-triple-overlap (source0.tex lines 1785--1792). The article's complete original proof of it is delivered with it (source0.tex lines 1794--1906, one \texttt{proof} environment of 113 lines). No mathematics is written, repaired or re-derived: the statement and the proof are the article's own lines, in the article's own notation, and no retained line is substituted or re-flowed.  No further source-local context is needed: every cross-reference of every retained line resolves inside the two delivered spans.  Nothing was omitted from any retained span, so the package carries no omission record.  No retained line contains a citation, so the wrapper carries no bibliography and no bibliography entry.  No retained line contains a figure environment, an image-inclusion command or drawing code, so no image is delivered and the package declares no asset dependency.  Editorial formatting only, in addition: an AMS \texttt{amsart} preamble that restates the article's own declarations and macros used by the retained lines, namely the environments \texttt{corollary} on separate counters, together with the standard packages those lines need.  The retained environments print in this document as Corollary 1 in order of appearance, whereas the article prints the same items with the numbering of its own full document.  The complete native source of the article as distributed in the arXiv e-print for this exact version is bundled unchanged as \texttt{sources/210182/source0.tex}.
\begin{corollary}[Exact overlap via interlacing]\label{cor:exact-triple-overlap}
Let $J_1<\cdots<J_h$ be ordered intervals.  The family
\[
 \gamma(J_a)+\gamma(J_b)+\gamma(J_c),
 \qquad a\le b\le c,
\]
has overlap multiplicity $O(h)$.
\end{corollary}

\begin{proof}
The reader should first realize that if each $J$ consisted of just one point, this would trivially follow from basic strict convexity. Indeed, the choice of the point $J_a$ within a fiber would uniquely determine the other two points $J_b,J_c$. The proof for intervals relies critically on interlacing.

Fix a point $P=(P_1,P_2)\in\mathbb{R}^2$. A triple of intervals
$(J_a,J_b,J_c)$, with $a\le b\le c$, contributes to $P$ if
\[
    P\in\gamma(J_a)+\gamma(J_b)+\gamma(J_c).
\]
This means that there are parameters
\[
    x\in J_a,\qquad y\in J_b,\qquad z\in J_c
\]
such that
\[
    P=\gamma(x)+\gamma(y)+\gamma(z),
\]
or, equivalently,
\[
    x+y+z=P_1,\qquad x^2+y^2+z^2=P_2.
\]
Since the intervals are ordered, we may arrange that $x\le y\le z$,
sorting the parameters within any interval that occurs more than once.
We must show that only $O(h)$ distinct triples of intervals can
contribute to this fixed $P$.

First, we use the convention that a shared endpoint belongs only to the
interval on its right, so that every parameter belongs to a unique
interval. For each contributing triple of intervals, choose one triple
of parameters $(x,y,z)$ as above.

Two chosen triples with the same smallest parameter $x$ must coincide.
Indeed,
\[
    y+z=P_1-x,\qquad y^2+z^2=P_2-x^2
\]
determine the ordered pair $y\le z$ uniquely. Since each parameter
belongs to a unique interval, coinciding parameter triples come from
the same triple of intervals.

We can therefore list the distinct contributing triples as
\[
    (J_{a_1},J_{b_1},J_{c_1}),\ldots,
    (J_{a_m},J_{b_m},J_{c_m}),
\]
with chosen parameters
\[
    x_\nu\in J_{a_\nu},\qquad
    y_\nu\in J_{b_\nu},\qquad
    z_\nu\in J_{c_\nu},
\]
in such a way that
\[
    x_1<x_2<\cdots<x_m.
\]

Consider two successive triples. Their parameters satisfy
\[
\begin{aligned}
    x_\nu+y_\nu+z_\nu
    &=x_{\nu+1}+y_{\nu+1}+z_{\nu+1}=P_1,\\
    x_\nu^2+y_\nu^2+z_\nu^2
    &=x_{\nu+1}^2+y_{\nu+1}^2+z_{\nu+1}^2=P_2.
\end{aligned}
\]
Applying Lemma~5.5 with
\[
\begin{aligned}
    (x,y,z)&=(x_\nu,y_\nu,z_\nu),\\
    (x',y',z')&=(x_{\nu+1},y_{\nu+1},z_{\nu+1}),
\end{aligned}
\]
we obtain
\[
    x_\nu<x_{\nu+1}
    \le y_{\nu+1}\le y_\nu
    \le z_\nu\le z_{\nu+1}.
\]
Thus the smallest and largest parameters move to the right, while the
middle parameter moves to the left. Because $J_1,\ldots,J_h$ are
ordered, their interval indices satisfy
\[
    a_\nu\le a_{\nu+1},\qquad
    b_\nu\ge b_{\nu+1},\qquad
    c_\nu\le c_{\nu+1}.
\]

At least one of these three indices changes at every step, since the
contributing triples of intervals are distinct. Each index takes
values in $\{1,\ldots,h\}$ and moves in only one direction, so each can
change at most $h-1$ times. Consequently,
\[
    m-1\le 3(h-1),
    \qquad\text{and hence}\qquad
    m\le 3h-2.
\]

Finally, allowing a shared endpoint to belong to both adjacent
intervals changes the count by at most an absolute factor. Under the
convention above, a parameter assigned to $J_j$ could originally have
belonged to $J_j$ or $J_{j-1}$. Thus each contributing triple under the
unique-endpoint convention accounts for at most $2^3$ original triples.
The number of contributing triples is therefore still $O(h)$, uniformly
in $P$.









\end{proof}

\end{document}
